\odiseachapter{nolan}{El cine teme a Homero}\label{cap:nolan}

Para mi decimoquinto cumpleaños, mi padre me había prometido un «regalo que te va a gustar» y yo sospechaba que sería un libro, porque los regalos de mi padre siempre eran libros.

En lo que se refiere a esos extraños animales, parecidos a ladrillos, con tapas de pasta y cuerpo de papel, yo era un depredador voraz y omnívoro. Leía todo lo que caía en mis manos y lo hacía sin contemplaciones, a toda velocidad.
A los quince años campeaba por la biblioteca de mi padre como Atila por el decadente Imperio romano. Solo que aquella biblioteca no era decadente, sino la cueva del tesoro.

El regalo llegó justo a tiempo para llevarme el enorme volumen a la playa. Fui a buscar a mi primo Pedro y los dos pasamos un buen rato admirándolo como quien venera una Biblia, maravillándonos con aquel guerrero griego, lanza en mano, que ocupaba la portada, acariciando el lomo azul, leyendo el título con la solemnidad de quien sabe que, a partir del momento en que nos sumergiéramos en las páginas sagradas, dejaríamos de ser niños para convertirnos en adultos. Peor aún, en poetas.

El libro era la preciosa versión de la \emph{Ilíada} y la \emph{Odisea} que el Círculo de Lectores publicó en 1971. La traducción ---¡de 1910, nada menos!--- era de don Luis Segalá y Estalella.
El primer canto arrancaba así:

\begin{quote}
Háblame, Musa, de aquel varón de multiforme ingenio que, después de destruir la sacra ciudad de Troya, anduvo peregrinando larguísimo tiempo, vió las poblaciones y conoció las costumbres de muchos hombres y padeció en su ánimo gran número de trabajos en su navegación por el ponto, en cuanto procuraba salvar su vida y la vuelta de sus compañeros á la patria. Mas ni aun así pudo librarlos, como deseaba, y todos perecieron por sus propias locuras. ¡Insensatos! Comiéronse las vacas del Sol, hijo de Hiperión; el cual no permitió que les llegara el día del regreso. ¡Oh diosa, hija de Júpiter!: cuéntanos aunque no sea más que una parte de tales cosas.
\end{quote}

Si se lee el párrafo de un tirón, el riesgo de sofoco es grande. El texto sonaba grandilocuente y lo era, pero también emocionante y exótico. «Insensatos», le gritaba a mi primo, «comiéronse las vacas del Sol, hijo de Hiperión». Y él respondía en el mismo tono: «¡Oh diosa, hija de Júpiter!: cuéntanos aunque no sea más que una parte de tales cosas», y los dos nos sentíamos bardos y guerreros y amigos de los olímpicos.

Aquel verano también leí la \emph{Ilíada} y cuando llegué al último canto, había decidido que ya no quería ser ni Ulises ---con quien tenía cuentas pendientes, que ya iré desgranando--- ni siquiera Aquiles. Quería ser Héctor de Troya, el más noble de todos los héroes del mundo.

Y, sin embargo, había un pasaje que no conseguía digerir. En el canto XXII, los troyanos se han refugiado tras las murallas de Troya, después de que Aquiles haya hecho una escabechina entre sus huestes, dejando el río Escamandro obturado de tanto cadáver. Héctor, sin embargo, se ha quedado fuera de la muralla. Su orgullo le impide ponerse a salvo, ya que la noche anterior, envalentonado por las sucesivas victorias contra los aqueos ---y por haber dado muerte al pobre Patroclo---, había rechazado volver a la ciudad tal como proponía el prudente Polidamante y ordenado acampar frente a los muros. Al día siguiente, Aquiles, furioso, vuelve a la batalla, decidido a vengar en sangre troyana la muerte de su amigo, y a los troyanos no les queda otra que ponerse a salvo. Todos menos Héctor. Estaba claro que quedarse a esperar al divino Pélida era un suicidio, pero, pensaba yo, nobleza obliga.

Y, sin embargo, cuando Aquiles deja de perseguir a Apolo, regresa a las murallas y se le echa encima, a Héctor le entra la flojera. La escena es así (la traducción es mía):

\begin{versocita}
Pensaba así Héctor mientras esperaba,\\
Ares belicoso, venía Aquiles,\\
dios de la guerra, del brillante casco,\\
así el Pélida, de fresno la lanza,\\
que esgrimía, terrible sobre el hombro,\\
como relámpago brillaba el bronce,\\
que ceñía su cuerpo, igual que hoguera,\\
o el sol alzándose en la madrugada.
\stanzabreak
Cuando Héctor lo vio, temblores le asieron,\\
y no pudo seguir junto a las puertas,\\
y Aquiles, confiando en sus pies veloces,\\
como un halcón persigue a una paloma,\\
así, velozmente, le perseguía.
\end{versocita}

Homero escribe \gr{ἕλε τρόμος} (\emph{héle trómos}), «le entró el temblor», y \gr{φοβηθείς} (\emph{phobētheís}), «huyendo aterrado». No hay duda alguna. Héctor ha visto cómo el semidiós se le echa encima, no ha podido controlar el pánico y se ha dado a la fuga.

Cuando leí aquella escena no daba crédito a mis ojos. ¿Cómo era posible, me preguntaba, que al valeroso general, después de tomar la decisión de enfrentar al Pélida ---por honor, obviamente--- le entrara el canguelo y saliera por piernas? Llegué incluso a escribir una variante de la historia en la que Héctor no tenía miedo, sino que intentaba agotar a Aquiles haciéndole correr, aparentemente sin darse cuenta de que se las estaba viendo con el de los pies veloces. Me costó bastantes años y muchas lecturas reconciliarme con el hecho de que el troyano, por noble que fuera, era humano y, por tanto, podía sentir miedo, como sentiría cualquiera frente a un enemigo invencible y furioso. Y cuando finalmente acepté la humanidad del héroe cuyo nombre lleva mi hijo, comprendí también que Homero me lo había dejado claro desde el principio. El poeta se limita a describir la huida con una compasión enorme, utilizando el símil de la pesadilla, en el que ni el perseguidor consigue alcanzar a su presa ni el perseguido escapar del que le acosa. Ni juzga a Aquiles por su saña, ni a Héctor por su miedo.

¿Cómo se traslada eso a la narrativa contemporánea? Si escogemos el cine como medio de comunicación de masas, no faltan versiones de la \emph{Ilíada}. En ninguna, que yo sepa ---desde luego no en la «clásica» \emph{Troya} de Petersen (2004), en la que Eric Bana le planta cara virilmente a Brad Pitt---, aparece la escena de la huida y, en consecuencia, el personaje de Héctor se aleja del original y no para bien. El héroe de Petersen es mucho más cliché que el de Homero (de Aquiles, mejor ni hablamos) y ese cliché se transmite a millones de espectadores.

Puede alegarse, con razón, que llevar una historia al cine es traducirla a otro lenguaje y, por tanto, interpretarla. Es decir, \emph{traduttore, traditore}. Pero ni todas las traducciones son equivalentes, ni todas las traiciones igual de imperdonables. Cuando la versión de Petersen se salta el miedo de Héctor, está optando por un personaje plano, el guerrero de valentía casi mecánica representado hasta la saciedad desde las novelas de caballerías. Para ese viaje no se precisaban tales alforjas.

Interpretar la \emph{Odisea} en clave cinematográfica requiere no poca osadía, más aún teniendo en cuenta que muchos otros cineastas lo han intentado antes. Y quizás ese sea mi mayor problema con la versión de Nolan. A pesar de su bardo rapero, a pesar de la habilidad con que quita de en medio a los dioses ---excepto a una onírica Atenea---, a pesar de elipsis muy inteligentes, como la que nos permite escuchar, pero no ver, a las sirenas, o de opciones arriesgadas y elegantes como escoger a Lupita Nyong'o para interpretar a Helena,
a pesar de que la película es entretenida y bellamente filmada, creo que la \emph{Nodisea} nos da gato por liebre.

Y es que Nolan no quiere limitarse a contarnos una de piratas. El argumento central de su película es que el asalto a Troya mediante la treta del caballo violó de forma imperdonable la ley de Zeus ---la \gr{ξενία} (\emph{xenía}), esto es, la obligación sagrada de dar hospitalidad al forastero--- y precipitó el fin de la Edad del Bronce. Pero esa tesis se sostiene a duras penas. Como veremos, el poema ofrece numerosos ejemplos que la contradicen y Nolan los atenúa o los elimina, forzando con ello una interpretación que no puede deducirse de los versos de la \emph{Odisea}.

La técnica de dar gato por liebre se repite constantemente en el film. Nolan transforma a una semidiosa (Helena, capítulo~\ref{cap:helena}) en una mujercita maltratada por su esposo, a una altiva divinidad (Circe, capítulo~\ref{cap:circe}) en una granjera, a un heredero valeroso, cruel y astuto (Telémaco, capítulo~\ref{cap:ninas}) en un jovencito confuso y al héroe más complicado de la literatura en un veterano con síndrome postraumático.

El problema es que, en los tiempos que corren, cuando el número de lectoras cada día se reduce más y el de lectores tiende a cero, uno se arriesga a que la \emph{Odisea} de Nolan o la \emph{Troya} de Petersen se conviertan en el canon. Y eso sería grave. Al salir del cine, uno tiene la sensación, como escribía Borges, de que \emph{hay algo que se queda}, sí: parte de la historia, una sombra de los personajes, una evocación más o menos afortunada de un tiempo y un lugar. Pero también \emph{hay algo que se queja}, y esa honda queja puede resumirse en una frase. El cine, o al menos el cine de masas, el cine de Hollywood, el cine de Nolan, le tiene miedo a Homero.
